\subsection{\texorpdfstring{以 \(\mu\) 为基的向量测度}{以 mu 为基的向量测度}}

定义 4. --- 设 \(\mu\) 是 T 上的正测度。称 T 上取值于 F 的向量测度 m 为以 \(\mu\) 为基的测度，如果存在 T 到 F 中的映射 f ，标量局部 \(\mu\)-可积，使得 \(\mathbf{m}(g) = \int g\mathbf{f}d\mu\) （对每个函数 \(g \in \mathcal{K}(\mathrm{T})\) ）。这时称 f 为 m 关于 \(\mu\) 的密度，并记 \(\mathbf{m} = \mathbf{f} \cdot \mu\) 。

立即看出，若 \(f_{1}\) 与 \(f_{2}\) 是 m 关于 \(\mu\) 的两个密度，则 \(f_{1}-f_{2}\) 标量局部 \(\mu\)-可忽略（Ch. V, §5, No. 3, 命题 3 的推论 2）；回忆一般说来这并不蕴含 \(f_{1}-f_{2}\) 局部几乎处处为零（参见 §1, 习题 12 与第 1 目的注记 2）。

设 \(\mathbf{m}\) 是以 \(\mu\) 为基、以 \(\mathbf{f}\) 为密度的测度。为使数值函数 \(g\) 关于 \(\mathbf{m}\) 本质可积，必须且只需 \(gf\) 标量本质 \(\mu\)-可积（Ch. V, §5, No. 3, 定理 1）。

命题 7. --- 设 f 是关于 T 上正测度 \(\mu\) 标量局部可积的映射，使得对每个函数 \(g \in \mathcal{K}(T)\) 有 \(\int g\mathbf{f} d\mu \in F\) 。则映射 \(g \mapsto \int g\mathbf{f} d\mu\) （ \(\mathcal{K}(T)\) 到 F 中的）是 T 上的向量测度，以 \(\mu\) 为基、以 f 为关于 \(\mu\) 的密度。

因为（第 1 目注记 3），只需证明：令 \(\mathbf{m}(g) = \int g\mathbf{f}d\mu\) ，则 \(\mathbf{z}' \circ \mathbf{m}\) 对每个 \(\mathbf{z}' \in \mathrm{F}'\) 是标量测度。但由于 \(\mathbf{z}'(\mathbf{m}(g)) = \int g\langle \mathbf{z}', \mathbf{f} \rangle d\mu\) ，有 \(\mathbf{z}' \circ \mathbf{m} = \langle \mathbf{z}', \mathbf{f} \rangle \cdot \mu\) ，由此即得我们的断言。

命题 8. --- 设 \(\mu\) 是 T 上的正测度， m 是 T 上取值于 F 的测度，以 \(\mu\) 为基、以 f 为关于 \(\mu\) 的密度。设 q 是 F 上的下半连续半范数。

\begin{enumerate}
\def\labelenumi{\alph{enumi})}
\item
  若对每个紧子集 \(\mathrm{K}\) （ \(\mathrm{T}\) 的），上积分 \(\int_{\mathrm{K}}^{*}(q \circ \mathbf{f}) d\mu\) 有限，则 \(\mathbf{m}\) 是 \(q\)-可优超的。
\item
  若 \(\mathbf{m}\) 是 \(q\)-可优超的，则 \(q(\mathbf{m})\) 以 \(\mu\) 为基；若进而 \(\mathbf{f}\) 是 \(\mu\)-可测的（作为 \(\mathrm{T}\) 到 \(\mathrm{F}_{\sigma}\) 中的映射），则 \(q \circ \mathbf{f}\) 局部 \(\mu\)-可积且 \(q(\mathbf{m}) = (q \circ \mathbf{f}) \cdot \mu\) 。
\end{enumerate}

\begin{enumerate}
\def\labelenumi{\alph{enumi})}
\tightlist
\item
  对每个有限子集 \(J\) （ \(A_q'\) 的），用 \(\lambda_J\) 表示诸测度 \(|\mathbf{z}' \circ \mathbf{m}|\) （其中 \(\mathbf{z}'\) 跑遍 \(J\) ）的上确界；若 \(g_J = \sup_{\mathbf{z}' \in J} |\langle \mathbf{z}', \mathbf{f} \rangle|\) ，则 \(\lambda_{J}=g_{J}\cdot\mu\) （Ch. V, §5, No. 2, 命题 2）。对 T 的每个相对紧开子集 U ，令 \(\lambda_{J,U}\) 为 \(\lambda_{J}\) 在 U 上的限制；首先证明：当 J 跑遍 \(A_{q}^{\prime}\) 的有限子集组成的有向集 F 时，族 \((\lambda_{J,U})\) 在 \(\mathcal{M}(U)\) 中有上界。的确，对每个函数 \(h\geqslant0\) （ \(K(U)\) 中的），
\end{enumerate}

\[
\int h d \lambda_ {\mathrm{J}, \mathrm{U}} = \int h g _ {\mathrm{J}} d \mu \leqslant \int^ {*} (q \circ \mathbf {f}) h d \mu \leqslant \| h \| \int_ {\mathrm{U}} ^ {*} (q \circ \mathbf {f}) d \mu ,
\]

由此即得我们的断言（Ch. II, §2, No. 2）。令 \(\nu_{\mathrm{U}}\) 为这族测度在 \(\mathcal{M}(\mathrm{U})\) 中的上确界。若 \(\mathrm{U}'\) 是 T 的另一个相对紧开子集，满足 \(\mathrm{U} \subset \mathrm{U}'\) ，则 \(\nu_{\mathrm{U}}\) 是 \(\nu_{\mathrm{U}'}\) 在 U 上的限制，这立即由递增有向测度集的上确界的表达式（Ch. II, §2, No. 2）以及 \(\lambda_{\mathrm{J,U}}\) 是 \(\lambda_{\mathrm{J,U}'}\) 到 U 上的限制这一事实得出。于是存在唯一的正测度 \(\nu\) ，它在每个 U 上的限制是 \(\nu_{\mathrm{U}}\) （Ch. III, §2, No. 1, 命题 1），且显然 \(\nu = q(\mathbf{m})\) 。

\begin{enumerate}
\def\labelenumi{\alph{enumi})}
\setcounter{enumi}{1}
\tightlist
\item
  由于诸测度 \(\lambda_{\mathrm{J}}\) 以 \(\mu\) 为基，其上确界 \(q(\mathbf{m})\) 亦然（Ch. V, §5, No. 5, 定理 2）。若 \(\mathbf{f}\) 是 \(\mu\)-可测的（对拓扑 \(\sigma(\mathrm{F},\mathrm{F}')\) ， \(\mathrm{F}\) 上的），则由定义立即得出，映射 \(g:t\mapsto(g_{\mathrm{J}}(t))_{\mathrm{J}\in\mathfrak{F}}\) （ \(\mathrm{T}\) 到乘积空间 \(\mathbf{R}^{\mathfrak{F}}\) 中的）是 \(\mu\)-可测的。 \(q\circ\mathbf{f}=\sup_{\mathrm{J}\in\mathfrak{F}}g_{\mathrm{J}}\) 到 \(\mathrm{T}\) 的每个使 \(g\) 连续的紧子集上的限制是下半连续的；从而 \(q \circ f\) 是 \(\mu\)-可测的（Ch. IV, §5, No. 5, 命题 8 的推论及 No. 10, 命题 15）。设 \(K\) 是 T 的紧子集；它容许一个分割，由一个 \(\mu\)-可忽略集与 \((\mathrm{K}_n)\) （ \(g\) 在其上连续的紧集序列）组成。这时 \(\int_{\mathrm{K}_n}^*\left(q\circ \mathbf{f}\right)d\mu = \sup_{\mathrm{J}}\int_{\mathrm{K}_n}g_{\mathrm{J}}d\mu \leqslant \int_{\mathrm{K}_n}dq(\mathbf{m})\) （对所有 \(n\) ）（Ch. IV, §1, No. 1, 定理 1 与 Ch. V, §7, No. 1, 命题 1），由此 \(\int_{\mathrm{K}}^{*}(q\circ \mathbf{f})d\mu = \sum_{n}\int_{\mathrm{K}_n}^*\left(q\circ \mathbf{f}\right)d\mu \leqslant \int_{\mathrm{K}}dq(\mathbf{m})\) 。但这证明了 \(q\circ \mathbf{f}\) 局部 \(\mu\)-可积，且 \(\lambda_{\mathrm{J}}\leqslant (q\circ \mathbf{f})\cdot \mu \leqslant q(\mathbf{m})\) （对所有 \(\mathbf{J}\in \mathfrak{F}\) ）；于是按定义， \(q(\mathbf{m}) = (q\circ \mathbf{f})\cdot \mu\) 。
\end{enumerate}

注记. --- 假设 \(A_{q}^{\prime}\) 中存在对 \(\sigma(\mathbf{F}^{\prime},\mathbf{F})\) 稠密的可数子集 D ；那么函数 \(q\circ f\) 总是 \(\mu\)-可测的，因为 \(q(\mathbf{f}(t)) = \sup_{\mathbf{z}^{\prime}\in\mathrm{D}} |\langle\mathbf{z}^{\prime},\mathbf{f}(t)\rangle|\) （Ch. IV, §5, No. 4, 定理 2 的推论 1）。于是，对

T 的每个紧子集 K ，有 \(\int_{K}^{*}(q\circ\mathbf{f})d\mu=\sup_{J}\int_{K}g_{J}d\mu\) ，其中 J 跑遍 D 的有限子集组成的可数有向集（Ch. IV, §1, No. 1, 定理 3 的推论）；于是可以看出，在这种情形，条件 \(\int_{K}^{*}(q\circ\mathbf{f})d\mu<+\infty\) （对 T 的每个紧子集 K ）是 m 为 q-可优超的必要充分条件。

命题 9. --- 设 F 是有限维 Banach 空间。T 上取值于 F 的每个测度 m 都是以 \textbar m\textbar 为基的测度。若 \(m = f \cdot |m|\) ，则 \(|f(t)| = 1\) 局部几乎处处成立（关于 \textbar m\textbar ）。为使 m 以 \(\mu\) 为基，必须且只需 \textbar m\textbar{} 以 \(\mu\) 为基，且若 \(m = g \cdot \mu\) 则 \(|m| = |g| \cdot \mu\) 。

设 \((\mathbf{e}_i)_{1 \leqslant i \leqslant n}\) 与 \((\mathbf{e}_i')_{1 \leqslant i \leqslant n}\) 是 F 与 \(\mathbf{F}'\) 的对偶基（A, II, §2, No. 6），满足 \(|\mathbf{e}_i'| = 1\) （对所有 \(i\) ）。这时 \(|\mathbf{e}_i' \circ \mathbf{m}| \leqslant |\mathbf{m}|\) （对每个指标 \(i\) ），因此（Ch. V, §5, No. 5, 定理 2） \(\mathbf{e}_i' \circ \mathbf{m} = h_i \cdot |\mathbf{m}|\) ，其中 \(h_i\) 有界且 \(|\mathbf{m}|\)-可测。令 \(\mathbf{h} = \sum_{i=1}^{n} h_i \cdot \mathbf{e}_i\) ，我们便有 \(\mathbf{m} = \mathbf{h} \cdot |\mathbf{m}|\) 。若 \(\mathbf{m} = \mathbf{f} \cdot |\mathbf{m}|\) ，命题 8 表明 \(|\mathbf{m}| = |\mathbf{f}| \cdot |\mathbf{m}|\) ，从而 \(|\mathbf{f}(t)| = 1\) 关于 \(|\mathbf{m}|\) 局部几乎处处成立（Ch. V, §5, No. 3, 命题 3 的推论 2）。最后的断言立即由命题 8 得出。

注记. --- 若 \(z = \sum_{i=1}^{n} z_i e_i\) ，则 \(\psi(z_1, \ldots, z_n) = |z|\) 是 \(R to the n\) 上的正齐次连续函数。令 \(\mu_i = e_i' \circ m = h_i \cdot |m|\) ，按定义（Ch. V, §5, No. 9）有 \(\psi(\mu_1, \ldots, \mu_n) = \psi(h_1, \ldots, h_n) \cdot |m| = |h| \cdot |m| = |m|\) 。

